Le gain de NvPD dépend de la divergence des séquences et de la façon dont le calcul est réparti. À un million de bases et 10\,\% d’éditions introduites, NvPD sur GPU devance les deux bibliothèques sur un cœur Grace pour les cinq paires testées (figure~\ref{fig:divergence-paires}). Pour des séquences proches, limiter les positions explorées réduit le travail. Sur quatre comparaisons de chromosomes HG002, le solveur associant coupes vérifiées et fronts bornés prend de 3,03 à 4,05 fois moins de temps que WFA2 sur un cœur EPYC (figure~\ref{fig:hg002-dispatch}). À l’échelle d’une collection, le partage entre paires évite les échanges à chaque ligne. Son efficacité forte atteint 72,70\,\% sur deux nœuds pour un lot synthétique, contre 4,59\,\% pour une paire carrée (figure~\ref{fig:mpifort}). Les arbres de la cohorte virale conservent leur topologie malgré deux scores surestimés (figure~\ref{fig:sensibilite-arbres}). À 100\,000 bases et 10\,\%, le classement entre NvPD et WFA2 change selon la paire. Des mesures à des taux intermédiaires permettraient de préciser la zone où leurs temps se rejoignent.
